% =============================================================================
% Presentación del proyecto final de NLP: sistema multi-agente ReAct (LangGraph) contra Naive RAG.
% Plantilla y estilo: la presentación de Computer Vision (turtlebot67/presentacion), con la portada
% slate-blocks y el tema minimalista en estilo/.
%
%   python presentacion/scripts/portada_fondo.py && (cd presentacion/figuras/gen && tectonic portada_fondo.tex)
%   cd presentacion && tectonic -X compile presentacion.tex     (o latexmk -pdf: dos pasadas, la portada usa remember picture)
%
% 10 diapositivas, una idea por diapositiva; el guion está en guion.md.
% Todos los números salen de results/: la corrida con Qwen2.5-7B y el juez de NVIDIA
% (results/qwen2.5-7b-instruct/, 26/11/2026), la de Mistral-7B (results/evaluacion.jsonl), y la
% validación del Auditor de GSM8K (results/validacion_gsm8k/). Las trazas son reales (evaluacion.jsonl),
% abreviadas donde dice «…».
% =============================================================================
\documentclass[aspectratio=1610, 11pt]{beamer}
% =============================================================================
% tema.tex — aspecto de todas las diapositivas salvo la portada.
%
% Tomado de la presentación de Computer Vision (turtlebot67/presentacion/estilo/tema.tex): fondo blanco,
% título arriba a la izquierda sin barra, pie con una franja gruesa (morado) y otra fina (celeste),
% navegación por secciones con un punto por diapositiva y número de diapositiva a la derecha.
% Texto en gris oscuro (cInk) y gris (cSrc), un solo acento (cAcc) para títulos y cifras clave.
%
% Cambios respecto al original: sin los estilos del robot y del circuito; colores para los dos
% métodos (multi-agente en morado, Naive RAG en gris), una caja de traza en monoespaciado para el
% bucle ReAct y los nodos del grafo de agentes.
% =============================================================================
\usepackage[T1]{fontenc}
\usepackage[utf8]{inputenc}
\usepackage[sfdefault,lining]{FiraSans}
\usepackage[scale=0.86]{FiraMono}
\usepackage[spanish,es-nodecimaldot,es-noshorthands]{babel}
\usepackage{graphicx}
\usepackage{amsmath}
\usepackage{booktabs}
\usepackage{array}
\usepackage{fontawesome5}
\usepackage{tikz}
\usepackage{pgfplots}
\pgfplotsset{compat=1.18}
\usetikzlibrary{positioning,arrows.meta,calc,fit,backgrounds,shapes.misc,shapes.geometric,babel}

% ---- paleta (la del original) -------------------------------------------------
\definecolor{cPub}{HTML}{8A97EE}     % morado
\definecolor{cAcc}{HTML}{4F5FC8}     % morado oscuro (títulos, acentos)
\definecolor{cLight}{HTML}{00BFFF}   % celeste
\definecolor{cCelD}{HTML}{0090C8}    % celeste oscuro (texto sobre blanco)
\definecolor{cEncBg}{HTML}{DEE2FB}   % morado claro
\definecolor{cDecBg}{HTML}{D9F4FF}   % celeste claro
\definecolor{cInk}{HTML}{2E3140}
\definecolor{cOurs}{HTML}{D92626}
\definecolor{cSrc}{HTML}{5A5D6E}
\definecolor{cMuted}{HTML}{8A8A96}
\definecolor{cRule}{HTML}{E2E3EA}
\definecolor{cPaper}{HTML}{F5F6FA}
\definecolor{cOk}{HTML}{3C9A5F}
% los dos métodos, iguales en todas las figuras
\colorlet{cMulti}{cAcc}
\colorlet{cNaive}{cMuted}

% ---- estructura de la plantilla -------------------------------------------------
\useoutertheme[subsection=false]{miniframes}
\useinnertheme{rectangles}
\setbeamertemplate{navigation symbols}{}
\setbeamersize{text margin left=10mm, text margin right=10mm}

\setbeamercolor{normal text}{fg=cInk, bg=white}
\setbeamercolor{background canvas}{bg=white}
\setbeamercolor{structure}{fg=cAcc}
\setbeamercolor{alerted text}{fg=cOurs}
\setbeamercolor{frametitle}{fg=cAcc, bg=white}
\setbeamercolor{footline1}{bg=cPub}
\setbeamercolor{footline2}{bg=cLight}
\setbeamercolor{section in head/foot}{fg=cAcc, bg=white}
\setbeamercolor{mini frame}{fg=cAcc}
\setbeamerfont{frametitle}{size=\Large, series=\mdseries}
\setbeamerfont{section in head/foot}{size=\fontsize{6.5}{7.5}\selectfont}
\setbeamertemplate{itemize item}{\color{cAcc}\raisebox{0.2ex}{\tiny$\blacksquare$}}
\setbeamertemplate{itemize subitem}{\color{cMuted}\raisebox{0.2ex}{\tiny$\blacktriangleright$}}
\setlength{\leftmargini}{1.1em}

\setbeamertemplate{headline}{}

\setbeamertemplate{frametitle}{%
  \vskip 4mm
  {\usebeamerfont{frametitle}\usebeamercolor[fg]{frametitle}\insertframetitle\par}%
  \vskip 0.8mm
}

\makeatletter
\setbeamertemplate{footline}{%
  \begin{beamercolorbox}[wd=\paperwidth, ht=2.2mm, dp=0pt]{footline1}\end{beamercolorbox}%
  \begin{beamercolorbox}[wd=\paperwidth, ht=0.6mm, dp=0pt]{footline2}\end{beamercolorbox}%
  \begin{beamercolorbox}[wd=\paperwidth]{section in head/foot}%
    \vskip 2pt
    \hbox to \paperwidth{\insertnavigation{\dimexpr\paperwidth-14mm\relax}\hfill
      \raisebox{0.4mm}{\fontsize{6}{7}\selectfont\color{cSrc}\insertframenumber\,/\,\inserttotalframenumber}\hskip 5mm}%
    \vskip 3pt
  \end{beamercolorbox}}
\makeatother

% ---- piezas reutilizables ------------------------------------------------------------------
% cifra grande con leyenda: \cifra{numero}{leyenda}
\newcommand{\cifra}[2]{{\color{cAcc}\fontsize{17}{19}\selectfont #1}\par\vspace{0.3mm}%
  {\scriptsize\color{cSrc}#2\par}}
% texto secundario
\newcommand{\nota}[1]{{\color{cSrc}#1}}
% pie de figura (un párrafo chico) y encabezado de una columna
\newcommand{\etq}[1]{{\fontsize{6.3}{7.6}\selectfont\color{cSrc}#1\par}}
\newcommand{\cab}[1]{{\footnotesize #1\par}\vspace{0.8mm}}
% los dos métodos en el texto
\newcommand{\multi}{{\color{cMulti}\textbf{multi-agente}}}
\newcommand{\naive}{{\color{cNaive}\textbf{Naive RAG}}}

\tikzset{
  % --- diagramas de bloques
  caja/.style={draw=cInk!45, line width=0.5pt, fill=white, rounded corners=2pt, align=center,
    inner sep=1.8pt, font=\fontsize{6.8}{8}\selectfont, text=cInk, execute at begin node={\hyphenpenalty=10000}},
  grupo/.style={draw=cInk!25, line width=0.5pt, fill=cPaper, rounded corners=3pt},
  ar/.style={-{Stealth[length=1.6mm]}, line width=0.6pt, draw=cInk!55},
  arAcc/.style={-{Stealth[length=1.6mm]}, line width=0.8pt, draw=cAcc},
  lbl/.style={font=\fontsize{6}{7}\selectfont, text=cSrc, align=center},
  % --- agentes del grafo: \node[agente] {...}
  agente/.style={draw=cAcc, line width=0.9pt, fill=cEncBg, rounded corners=3pt, align=center,
    inner xsep=3pt, inner ysep=2.5pt, font=\footnotesize, text=cInk},
  herramienta/.style={draw=cCelD, line width=0.6pt, fill=cDecBg, rounded corners=2pt, align=center,
    inner sep=2pt, font=\fontsize{6.8}{8}\selectfont, text=cInk},
  fin/.style={draw=cInk!45, line width=0.5pt, fill=white, rounded corners=6pt, align=center,
    inner sep=2.2pt, font=\fontsize{6.8}{8}\selectfont, text=cInk},
  % --- traza ReAct en monoespaciado: \node[traza] {...}
  traza/.style={draw=cRule, line width=0.5pt, fill=cPaper, rounded corners=2pt, align=left,
    inner sep=2.2mm, font=\ttfamily\fontsize{5.6}{6.9}\selectfont, text=cInk},
}
% palabras clave de una traza
\newcommand{\kw}[1]{{\color{cAcc}\bfseries #1}}
\newcommand{\obs}[1]{{\color{cCelD}#1}}

% ---- pgfplots: ejes sobrios, en la fuente y los colores de la presentación
\pgfplotsset{
  sobrio/.style={
    axis line style={draw=cInk!45}, tick style={draw=cInk!45},
    tick label style={font=\fontsize{6}{7}\selectfont, text=cSrc},
    label style={font=\fontsize{6.5}{7.5}\selectfont, text=cSrc},
    legend style={font=\fontsize{6}{7}\selectfont, draw=none, fill=none, text=cInk},
    grid style={draw=cRule}, major grid style={draw=cRule},
    every axis plot/.append style={line width=0.9pt},
  },
}

% =============================================================================
% portada-slate.tex — SOLO la portada, adaptada de "slate-blocks" (style.tex).
%
% Original: slate-blocks, a Beamer theme with seven named blocks, (c) its
% contributors, CC BY 4.0 (https://creativecommons.org/licenses/by/4.0/);
% licencia completa en estilo/LICENSE-slate-blocks.
%
% Tomada de la presentación de Computer Vision (turtlebot67/presentacion/estilo/portada-slate.tex).
% Cambios respecto al original:
%   * se toman solo la paleta, \slatecover y \slatetitlepage; NO se cargan
%     Madrid/miniframes/circles, ni los colores de beamer, ni los bloques, ni la
%     marca de agua: el resto de la presentación usa estilo/tema.tex;
%   * fondo: una traza ReAct real de la corrida (figuras/gen/portada_fondo.pdf, generada desde
%     figuras/portada_fondo.tex) en lugar de la fotografía del tema;
%   * logotipo más pequeño (abajo a la derecha);
%   * autores y curso en un bloque con velo blanco para que se lean sobre el fondo.
% =============================================================================

% ---- la paleta del tema (un color y sus derivados) --------------------------
\definecolor{ThemeColor}{RGB}{190, 194, 203}              % #bec2cb
\definecolor{AnalogousColor-1}{RGB}{193, 190, 203}        % #c1becb
\definecolor{AnalogousColor-2}{RGB}{190, 201, 203}        % #bec9cb
\definecolor{TetradicColor-1}{RGB}{203, 190, 194}         % #cbbec2
\definecolor{TetradicColor-2}{RGB}{203, 199, 190}         % #cbc7be
\definecolor{TetradicColor-3}{RGB}{190, 203, 199}         % #becbc7
\definecolor{SplitComplementaryColor-1}{RGB}{203, 193, 190} % #cbc1be
\definecolor{SplitComplementaryColor-2}{RGB}{201, 203, 190} % #c9cbbe

% ---- una imagen que cubre la pagina sin deformarse ---------------------------
\makeatletter
\newsavebox{\slate@coverbox}
\newcommand{\slatecover}[1]{%
  \sbox{\slate@coverbox}{\includegraphics[width=\paperwidth]{#1}}%
  \ifdim\ht\slate@coverbox<\paperheight
    \includegraphics[height=\paperheight]{#1}%
  \else
    \usebox{\slate@coverbox}%
  \fi
}
\makeatother

% ---- campos de la portada ----------------------------------------------------
\newcommand{\slateBackground}{figuras/gen/portada_fondo.pdf}
\newcommand{\slateTitle}    {\large\textbf{\textrm{Title}}}
\newcommand{\slateSubTitle} {\small\textrm{Subtitle}}
\newcommand{\slateAuthor}   {Author}
\newcommand{\slateAdvisor}  {}
\newcommand{\slateAffiliate}{Affiliation}
\newcommand{\slateDate}     {\today}
\newcommand{\slateLogo}     {\includegraphics[width=0.17\paperwidth]{figuras/logo_utec.png}}

% ---- la portada ----------------------------------------------------------------
\newcommand{\slatetitlepage}{%
  \begin{tikzpicture}[remember picture, overlay]
    \node[inner sep=0pt, opacity=0.55] at (current page.center)
      {\slatecover{\slateBackground}};
    \node[
      above=0.2cm, align=center,
      draw=AnalogousColor-1!50,
      left color=AnalogousColor-1!75, right color=AnalogousColor-1!25,
      shading angle=60,
      double, double distance=0.1cm,
      inner xsep=15pt, inner ysep=10pt,
      minimum width=0.74\textwidth, text width=0.66\textwidth,
      color=AnalogousColor-2!250,
    ] (title) at (current page.center)
      {\LARGE\slateTitle\ifx\slateSubTitle\empty\else\\[5pt] \small\slateSubTitle\fi};   % sin subtitulo: \def\slateSubTitle{}
    % autores, curso, filiacion y fecha en un solo bloque con un velo blanco: sobre la imagen
    % el texto gris del tema original no se leia
    % solo los nombres, uno por linea, a la izquierda debajo del titulo (sin profesor, curso ni fecha)
    \node[below=0.4cm, anchor=north west, align=left, fill=white, fill opacity=0.78, text opacity=1,
      rounded corners=3pt, inner xsep=10pt, inner ysep=5pt, text=AnalogousColor-1!250]
      (author) at (title.south west)
      {\slateAuthor};
    \node[fill=white, fill opacity=0.85, text opacity=1, rounded corners=3pt, inner sep=4pt,
      anchor=south east] at ([xshift=-0.4cm, yshift=0.35cm]current page.south east) {\slateLogo};
  \end{tikzpicture}%
}


\title[Multi-agente ReAct]{Sistema multi-agente ReAct}
\author[Grupo]{Felix, Huamani y Quispe}
\institute{UTEC}
\date{2026}

\renewcommand{\slateTitle}{\large\textbf{\textrm{Sistema multi-agente ReAct:\\agentes que verifican antes de responder}}}
\def\slateSubTitle{}
\renewcommand{\slateAuthor}{\small Renzo Josimar Felix Aponte\\\small Jose Eduardo Huamani Ñaupas\\\small Diego Randolp Quispe Apaza\\[2pt]\scriptsize Procesamiento de Lenguaje Natural \textperiodcentered\ UTEC \textperiodcentered\ 2026}

% ---- piezas con argumentos (en el preámbulo: beamer lee el cuerpo de un frame como argumento)
% una lección: \leccion{cifra}{afirmación}
\newcommand{\leccion}[2]{%
  \begin{tabular}{@{}p{30mm}@{\hspace{3mm}}p{\dimexpr\textwidth-33mm\relax}@{}}
    \leavevmode{\color{cAcc}\fontsize{15}{17}\selectfont #1} &
    \hyphenpenalty=10000\raggedright\leavevmode{\small #2\par}
  \end{tabular}\par\vspace{2mm}}
% una fila de la tabla de herramientas: \tarea{icono}{tarea}{dataset}{herramientas}{auditoría}
\newcommand{\tarea}[5]{\textcolor{cAcc}{#1}\ \textbf{#2} & #3 & #4 & #5 \\}
% una métrica con su barra: \par de barras horizontales para los dos métodos
\pgfplotsset{barras/.style={sobrio, xbar, bar width=5pt, y dir=reverse, xmin=0, xmax=1.08,
    axis x line=none, axis y line*=left, y axis line style={draw=none}, ytick style={draw=none},
    ytick=data, yticklabel style={font=\fontsize{7}{8}\selectfont, text=cInk, align=right, text width=27mm},
    scale only axis, enlarge y limits={abs=0.55},
    nodes near coords, nodes near coords style={font=\fontsize{6.5}{7.5}\selectfont, text=cInk,
      /pgf/number format/.cd, fixed, precision=2, zerofill, assume math mode=true},
    legend style={at={(0.5,1.02)}, anchor=south, legend columns=2, /tikz/every even column/.append style={column sep=4mm}}}}
\newcommand{\si}{{\color{cOk}\faCheck}}
\newcommand{\no}{{\color{cOurs}\faTimes}}

\begin{document}

% =============================================================================
% 1 · PORTADA
% =============================================================================
{\setbeamertemplate{headline}{}\setbeamertemplate{footline}{}
\begin{frame}[plain]
  \slatetitlepage
\end{frame}}

% =============================================================================
\section{Sistema}
% =============================================================================

% 2 · el reto y la arquitectura -------------------------------------------------
\begin{frame}{El reto y nuestra arquitectura}
  % \noindent: un { justo después del título, beamer lo toma como subtítulo
  \noindent{\small Un LLM de 7B en 4 bits (\textbf{Qwen2.5-7B} en una Tesla T4) que \textbf{no entregue alucinaciones}:
  tres agentes en un \textbf{grafo cíclico de LangGraph}, con herramientas reales y un Auditor que puede rechazar.\par}
  \vspace{5mm}
  \centering
  \begin{tikzpicture}[x=0.83mm, y=1mm]
    \node[fin] (q) at (0,0) {pregunta};
    \node[herramienta] (kb) at (22,7) {\faDatabase\ \textbf{kb}\\ChromaDB};
    \node[herramienta] (web) at (22,-7) {\faGlobe\ \textbf{web}\\DuckDuckGo};
    \node[agente, text width=31mm] (inv) at (58,0) {\textbf{Investigador}\\[0.3mm]{\scriptsize Thought \textrightarrow\ Action \textrightarrow\ Observation}};
    \node[agente, text width=20mm] (aud) at (95,0) {\textbf{Auditor}\\[0.3mm]{\scriptsize de Hechos}};
    \node[agente, text width=15mm] (red) at (122,0) {\textbf{Redactor}};
    \node[fin] (out) at (143,0) {respuesta};
    \node[fin, text=cSrc] (abs) at (95,-19) {abstención};
    \node[herramienta, anchor=north] (tools) at (58,-11) {kb \textperiodcentered\ web \textperiodcentered\ calc \textperiodcentered\ python};
    \draw[ar] (q) -- (kb.west); \draw[ar] (q) -- (web.west);
    \draw[ar] (kb.east) -- (inv); \draw[ar] (web.east) -- (inv);
    \node[lbl] at (22,0) {en paralelo};
    \draw[ar, {Stealth[length=1.4mm]}-{Stealth[length=1.4mm]}, draw=cCelD] (inv) -- (tools);
    \draw[arAcc] (inv) -- (aud);
    \draw[arAcc] (aud) -- node[above, font=\scriptsize, text=cOk]{\faCheck} (red);
    \draw[ar] (red) -- (out);
    \draw[ar, draw=cOurs!80] (aud.north) to[out=135, in=45] node[lbl, above, text=cOurs!80!black]{rechazo + motivo} (inv.north);
    \draw[ar] (aud) -- node[lbl, right]{3 rechazos} (abs);
  \end{tikzpicture}
  \vspace{5mm}

  \begin{columns}[T,onlytextwidth]
    \begin{column}{0.32\textwidth}\cifra{3 agentes}{con roles separados: buscar, verificar, redactar}\end{column}
    \begin{column}{0.32\textwidth}\cifra{4 herramientas}{base vectorial, web, calculadora y Python aislado}\end{column}
    \begin{column}{0.32\textwidth}\cifra{abstenerse}{antes que entregar una respuesta sin respaldo}\end{column}
  \end{columns}
\end{frame}

% 3 · el bucle ReAct con un rechazo real ---------------------------------------------
\begin{frame}{El bucle ReAct y un rechazo real}
  \centering
  \begin{tikzpicture}[x=1mm, y=1mm, every node/.append style={anchor=north west}]
    \node[traza, text width=37mm] (q) at (0,0) {%
      \textsf{\footnotesize\bfseries MBPP 11}\\[1mm]
      Write a python function to remove first\\ and last occurrence of a given character.\\[1mm]
      assert remove\_Occ("hello","l") == "heo"\\
      assert remove\_Occ("abcda","a") == "bcd"\\
      assert remove\_Occ("PHP","P") == "H"};
    \node[traza, text width=37mm, draw=cOurs!60, fill=cOurs!5] (r1) at (0,-31) {%
      \textsf{\footnotesize\bfseries Ronda 1 \textperiodcentered\ \no\ rechazada}\\[1mm]
      \textsf{\scriptsize El Auditor \textbf{ejecuta} el código con las pruebas:}\\[0.8mm]
      assert remove\_Occ("abcda","a") == "bcd"\\
      {\color{cOurs}AssertionError}};
    \node[traza, text width=80mm] (r2) at (45,0) {%
      \textsf{\footnotesize\bfseries Ronda 2 \textperiodcentered\ el error vuelve al Investigador}\\[1mm]
      \kw{Thought:} The previous solution did not work because it\\ used `replace' \dots{} We need to find the indices of the\\ first and last occurrences \dots\\
      \kw{Action:} python[def remove\_Occ(s, ch):\\
      \ \ \ \ first\_index = s.find(ch)\\
      \ \ \ \ last\_index = s.rfind(ch)\\
      \ \ \ \ \dots\\
      \ \ \ \ return s[:first\_index] + s[first\_index+1:last\_index]\\
      \ \ \ \ \ \ \ \ \ \ \ + s[last\_index+1:]\\[0.6mm]
      assert remove\_Occ("abcda", "a") == "bcd" \dots\\
      print("ok")]\\
      \kw{Observation:} \obs{ok}\\
      \kw{Thought:} The code passed all the tests \dots\\
      \kw{Final:} \textsf{el mismo código}};
    \node[traza, text width=80mm, draw=cOk!70, fill=cOk!6] (ok) at ([yshift=-2mm]r2.south west) {%
      \textsf{\footnotesize\bfseries Auditor \textperiodcentered\ \si\ aprobada:} \textsf{\scriptsize las 3 pruebas pasan; el Redactor entrega el código sin reescribirlo}};
    \draw[ar, draw=cOurs!70] (r1.east) -- ++(4,0) |- ([yshift=-8mm]r2.north west);
  \end{tikzpicture}\\[1mm]
  {\scriptsize\color{cSrc}Traza real de la corrida, abreviada en «…». El modelo escribe Thought y Action; la Observation la escribe el sistema.\par}
\end{frame}

% 4 · herramientas y datos ------------------------------------------------------------
\begin{frame}{Un grafo por tarea, con su herramienta y su verificación}
  \vspace{2mm}
  {\small
  \renewcommand{\arraystretch}{1.55}
  \begin{tabular}{@{}>{\raggedright\arraybackslash}m{25mm}>{\raggedright\arraybackslash}m{16mm}>{\raggedright\arraybackslash}m{34mm}>{\raggedright\arraybackslash}m{50mm}@{}}
    \toprule
    \textcolor{cSrc}{\footnotesize Tarea} & \textcolor{cSrc}{\footnotesize Dataset} & \textcolor{cSrc}{\footnotesize Herramientas} & \textcolor{cSrc}{\footnotesize Cómo verifica el Auditor} \\
    \midrule
    \tarea{\faQuestionCircle}{Preguntas}{WikiQA}{\texttt{kb} + \texttt{web}, en paralelo al inicio}{el LLM contrasta la respuesta con la evidencia}
    \tarea{\faCalculator}{Matemáticas}{GSM8K}{\texttt{calc}, calculadora segura}{resuelve el problema por su cuenta y compara el número}
    \tarea{\faCode}{Código}{MBPP}{\texttt{python}, proceso aislado}{\textbf{ejecuta} el código con los \texttt{assert} del enunciado}
    \bottomrule
  \end{tabular}\par}
  \vspace{7mm}
  \begin{columns}[T,onlytextwidth]
    \begin{column}{0.32\textwidth}\cifra{1119 frases}{de WikiQA en ChromaDB, también las incorrectas: el Naive RAG no recibe la respuesta servida}\end{column}
    \begin{column}{0.32\textwidth}\cifra{3 rondas}{máximo por consulta; 5 en MBPP, donde reintentar no deja pasar código erróneo}\end{column}
    \begin{column}{0.32\textwidth}\cifra{sin web}{en GSM8K y MBPP: podría devolver la solución publicada del benchmark}\end{column}
  \end{columns}
\end{frame}

% =============================================================================
\section{Evaluación}
% =============================================================================

% 5 · cómo evaluamos -------------------------------------------------------------------
\begin{frame}{Cómo evaluamos: 50 consultas, dos métodos, un juez}
  \begin{columns}[T,onlytextwidth]
    \begin{column}{0.5\textwidth}
      \cab{\textbf{Métricas} \nota{(definiciones de RAGAS, con LLM-as-a-Judge)}}
      \vspace{1mm}
      {\small\textbf{Faithfulness}\par}\vspace{1.2mm}
      \hspace*{5mm}$\dfrac{\text{\footnotesize afirmaciones respaldadas por el contexto}}{\text{\footnotesize afirmaciones de la respuesta}}$\par\vspace{3mm}
      {\small\textbf{Answer Relevance}\par}\vspace{1.2mm}
      \hspace*{5mm}$\dfrac{\text{\footnotesize puntaje del juez, de 0 a 10}}{\text{\footnotesize 10}}$\par\vspace{3mm}
      {\small\textbf{Verificación objetiva}: número final (GSM8K) y \texttt{assert} ejecutados (MBPP)\par}
      \vspace{2mm}
      {\footnotesize\color{cSrc}Una abstención puntúa 0 en las dos métricas del juez.\par}
    \end{column}
    \begin{column}{0.46\textwidth}
      \cab{\textbf{Mismas 50 consultas} \nota{20 WikiQA \textperiodcentered\ 15 GSM8K \textperiodcentered\ 15 MBPP}}
      \vspace{2mm}
      \begin{tikzpicture}[x=1mm, y=1mm]
        \node[agente, text width=52mm, anchor=north west] at (0,0) {\textbf{multi-agente}\\{\scriptsize herramientas, Auditor y hasta 3 rondas}};
        \node[caja, text width=52mm, anchor=north west, font=\footnotesize, minimum height=9mm] at (0,-13) {\textbf{Naive RAG}\\{\scriptsize 1 recuperación (3 fragmentos) y 1 llamada, sin auditoría}};
        \node[font=\scriptsize, text=cSrc, anchor=north west, text width=50mm] at (0,-26) {los dos con el mismo \textbf{Qwen2.5-7B} de 4 bits};
        \node[herramienta, text width=52mm, anchor=north west] at (0,-34) {\faBalanceScale\ \textbf{juez independiente}: Nemotron 120B (API de NVIDIA), el mismo para ambos};
      \end{tikzpicture}
    \end{column}
  \end{columns}
\end{frame}

% 6 · resultados globales -------------------------------------------------------------------
\begin{frame}{Resultados sobre las 50 consultas}
  \begin{columns}[c,onlytextwidth]
    \begin{column}{0.58\textwidth}\centering
      \begin{tikzpicture}
        \begin{axis}[barras, width=46mm, height=46mm, ytick={0,1,2,3},
            yticklabels={Faithfulness, Answer Relevance, {Exactitud\\{\scriptsize\color{cSrc}GSM8K y MBPP}}, {Incorrectas entregadas\\{\scriptsize\color{cSrc}menor es mejor}}}]
          \addplot[fill=cMulti, draw=none, bar shift=3pt] coordinates {(0.804,0) (0.860,1) (0.733,2) (0.100,3)};
          \addplot[fill=cNaive!70, draw=none, bar shift=-3pt] coordinates {(0.678,0) (0.960,1) (0.833,2) (0.167,3)};
          \legend{multi-agente, Naive RAG}
        \end{axis}
      \end{tikzpicture}
    \end{column}
    \begin{column}{0.4\textwidth}
      \cifra{0.914 vs 0.696}{Faithfulness en las 44 consultas que responden los dos}
      \vspace{3.5mm}
      \cifra{0.977 vs 0.973}{Answer Relevance en esas mismas 44}
      \vspace{3.5mm}
      \cifra{6 de 50}{abstenciones del multi-agente: puntúan 0 y bajan sus promedios}
    \end{column}
  \end{columns}
  \vspace{2mm}
  {\footnotesize\color{cSrc}Gana en fidelidad y entrega menos errores. Pierde en relevancia y exactitud por GSM8K (siguiente diapositiva).\par}
\end{frame}

% 7 · por dataset -------------------------------------------------------------------------
\begin{frame}{Dónde gana y dónde pierde}
  \vspace{2mm}
  \begin{columns}[T,onlytextwidth]
    \begin{column}{0.31\textwidth}
      {\color{cAcc}\faQuestionCircle}\ \textbf{WikiQA} {\color{cOk}\faCheck}\par\vspace{3mm}
      \cifra{0.84 vs 0.66}{Faithfulness: responde con lo que recuperó}
      \vspace{3mm}
      \cifra{0.91 vs 0.90}{Answer Relevance}
    \end{column}
    \begin{column}{0.31\textwidth}
      {\color{cAcc}\faCalculator}\ \textbf{GSM8K} {\color{cOurs}\faTimes}\par\vspace{3mm}
      \cifra{9 vs 14}{aciertos de 15}
      \vspace{3mm}
      \cifra{3}{abstenciones con la respuesta correcta en la mano: el Auditor la rechazó}
    \end{column}
    \begin{column}{0.31\textwidth}
      {\color{cAcc}\faCode}\ \textbf{MBPP} {\color{cOk}\faCheck}\par\vspace{3mm}
      \cifra{13 vs 11}{aciertos de 15}
      \vspace{3mm}
      \cifra{0 vs 4}{códigos incorrectos entregados}
    \end{column}
  \end{columns}
  \vspace{6mm}
  {\footnotesize\color{cSrc}Primero el multi-agente, después el Naive RAG. Qwen2.5-7B, juez Nemotron 120B.\par}
\end{frame}

% =============================================================================
\section{Hallazgos}
% =============================================================================

% 8 · el juez y el modelo importan ------------------------------------------------------------
\begin{frame}{El juez importa tanto como el modelo}
  \begin{columns}[c,onlytextwidth]
    \begin{column}{0.58\textwidth}
      {\small
      \renewcommand{\arraystretch}{1.35}
      \begin{tabular}{@{}l>{\centering\arraybackslash}p{20mm}>{\centering\arraybackslash}p{20mm}@{}}
        & \textbf{Mistral-7B} & \textbf{Qwen2.5-7B} \\
        & \textcolor{cSrc}{\footnotesize juez de 7B} & \textcolor{cSrc}{\footnotesize juez de 120B} \\
        \midrule
        \textcolor{cSrc}{\footnotesize multi / naive} & & \\
        Faithfulness & 0.58 / \textcolor{cAcc}{0.66} & \textcolor{cAcc}{0.80} / 0.68 \\
        Answer Relevance & 0.79 / \textcolor{cAcc}{0.91} & 0.86 / \textcolor{cAcc}{0.96} \\
        Exactitud & \textcolor{cAcc}{0.50} / 0.30 & 0.73 / \textcolor{cAcc}{0.83} \\
        Incorrectas $\downarrow$ & \textcolor{cAcc}{0.23} / 0.70 & \textcolor{cAcc}{0.10} / 0.17 \\
        \bottomrule
      \end{tabular}\par}
    \end{column}
    \begin{column}{0.38\textwidth}
      \cifra{0.86 > 0.73}{con Mistral, el juez de 7B daba más fidelidad a las respuestas \textbf{incorrectas} del multi-agente que a las correctas}
      \vspace{4mm}
      \cifra{0.36 < 0.78}{el juez de 120B sí las separa (respuestas del Naive RAG)}
    \end{column}
  \end{columns}
  \vspace{4mm}
  {\footnotesize\color{cSrc}El enunciado exige que el orquestador sea local; el juez no, así que lo hicimos independiente.\par}
\end{frame}

% 9 · un guardrail vale lo que su verificación --------------------------------------------------
\begin{frame}{Un guardrail vale lo que su verificación}
  \begin{columns}[T,onlytextwidth]
    \begin{column}{0.47\textwidth}
      {\color{cAcc}\faCode}\ \textbf{MBPP: el Auditor ejecuta}\par\vspace{2mm}
      \cifra{0}{códigos incorrectos entregados, con los dos modelos}
      \vspace{2mm}
      {\footnotesize En 3 de 15 consultas el error de ejecución volvió al Investigador: corrigió una y se abstuvo en las otras dos, donde el Naive RAG también falló.\par}
    \end{column}
    \begin{column}{0.49\textwidth}
      {\color{cAcc}\faCalculator}\ \textbf{GSM8K: el Auditor opinaba}\par\vspace{2mm}
      \begin{tikzpicture}[x=1mm, y=1mm]
        \node[traza, text width=56mm, anchor=north west] at (0,0) {%
          \textsf{\scriptsize Respuesta del Investigador: \textbf{3 bolts} (correcta)}\\[0.8mm]
          \kw{Auditor:} El cálculo no es correcto. \dots\ el total\\ es 2 + 1 = 3 bolts, pero el cálculo debe ser\\ 2 + 2/2 = 3, lo cual es correcto, pero \dots\\ {\color{cOurs}RECHAZADO} \textsf{\scriptsize (3 veces: abstención)}};
      \end{tikzpicture}\par\vspace{2mm}
      {\footnotesize \textbf{Cambio:} el Auditor resuelve el problema \textbf{sin ver la respuesta} y solo rechaza si los números no coinciden. Validado en 20 problemas que no están en la evaluación.\par}
    \end{column}
  \end{columns}
\end{frame}

% =============================================================================
\section{Cierre}
% =============================================================================

% 10 · lecciones + gracias ------------------------------------------------------------------
\begin{frame}{Lecciones aprendidas}
  \vspace{2mm}
  \leccion{0 de 15}{Verificar ejecutando, no opinando: ningún código erróneo pasó al Auditor de MBPP.}
  \leccion{0.36 vs 0.78}{Un juez débil no distingue una respuesta correcta de una incorrecta: hay que medir con uno independiente.}
  \leccion{6 de 50}{Abstenerse baja las métricas promediadas, pero evita entregar errores.}
  \leccion{9 de 20}{El 7B respondía en español a preguntas en inglés: hay que nombrarle el idioma.}
  \vfill
  {\fontsize{16}{19}\selectfont\color{cAcc}¡Gracias!}\ \ \nota{\small ¿Preguntas?}\par
  \vspace{1mm}
\end{frame}

\end{document}
